\chapter{Quine Structures and Self-Rewriting Atoms}
\label{ch:quines}

This chapter develops the self-referential core of Polyxan algebra.  
We introduce self-rewriting atoms, construct the category of
quine-enriched hypergraphs, define the quine operator 
$\mathcal{Q}:\mathbf{TypHyper}\to\mathbf{QuinePolyx}$,
prove the Quine Fixed-Point Theorem, define higher-order quines and their
stabilisation, establish the Polyxan Löb Theorem, and identify the universal
media quine as the initial algebra of the higher-quine monad.

This chapter plays the same role for Polyxan hypergraphs that the
semantic conformal point $\mathcal{L}^{\ast}$ plays for RSVP fields:  
the locus where self-consistency, symmetry, and curvature-entropy balance
force the system into a universal fixed point.

% ================================================================
\section{Self-Rewriting Atoms}
% ================================================================

\begin{definition}[Self-Rewriting Atom]
A content atom $a=(t,\mathrm{payload},\sigma,w)$ is a 
\emph{self-rewriting atom} if its payload encodes a rewrite rule
\[
\rho_a : L_a \longleftarrow K_a \longrightarrow R_a
\]
such that $R_a$ contains an isomorphic copy of $a$ as a sub-hypergraph.
\end{definition}

Thus $a$ performs a rewrite in which $a$ is reproduced.  
Self-rewriting atoms are the discrete analogues of reflexive operators,
Gödelian self-reference, and self-similar modes at the RSVP conformal point.

% ================================================================
\section{The Category \texorpdfstring{$\mathbf{QuinePolyx}$}{QuinePolyx}}
% ================================================================

\begin{definition}[Quine-Enriched Polyxan Hypergraph]
A \emph{quine-enriched hypergraph} is a pair $(G,A_{\mathrm{quine}})$ where  
$G$ is a typed Polyxan hypergraph and   
$A_{\mathrm{quine}} \subseteq V(G)$ is a set of self-rewriting atoms.
\end{definition}

A morphism $(G,A_{\mathrm{quine}})\to(H,B_{\mathrm{quine}})$ is a typed
hypergraph homomorphism preserving quine-atoms.

Let
\[
U:\mathbf{QuinePolyx}\to\mathbf{TypHyper}
\]
be the forgetful functor.

% ================================================================
\section{The Quine Operator \texorpdfstring{$\mathcal{Q}$}{Q}}
% ================================================================

\begin{definition}[Quine Operator]
The \emph{quine operator}
\[
\mathcal{Q} : \mathbf{TypHyper} \longrightarrow \mathbf{QuinePolyx}
\]
freely adjoins to each vertex $v\in V(G)$ a self-rewriting atom $q_v$ whose 
payload contains the EHR rewrite system that regenerates $G$.
\end{definition}

This yields the adjunction:
\[
\mathcal{Q} \;\dashv\; U.
\]

% ================================================================
\section{Iterated Quine Completion and the Fixed-Point Theorem}
% ================================================================

Define the tower:
\[
G_0=G,\quad G_{n+1}=\mathcal{Q}(G_n).
\]

\begin{definition}
The \emph{quine completion} of $G$ is
\[
G_\infty := \varinjlim_{n\to\infty} G_n.
\]
\end{definition}

\begin{theorem}[Quine Fixed-Point Theorem]
\label{thm:quine-fixed-point}
For any finite $G$, the colimit $G_\infty$ exists, is finite, and satisfies:
\[
\mathcal{Q}(G_\infty)\cong G_\infty.
\]
Moreover, $G_\infty$ is the unique media quine reachable from $G$ under EHR and
quine closure.
\end{theorem}

\begin{proof}[Sketch]
Each $\mathcal{Q}$-step strictly decreases the curvature functional $\mathcal{H}_0$
unless the needed quine structure is already present.  
Thus the tower stabilises.  
Fixed-pointness follows from self-reproduction of the quine atoms.
\end{proof}

% ================================================================
\section{Higher-Order Quines}
% ================================================================

\begin{definition}[Order-$k$ Quine]
An \emph{order-$k$ quine} $q^{(k)}$ is a self-rewriting atom whose payload 
rewrites the rewrite rules of order-$(k-1)$ quines.
\end{definition}

Thus:
\[
q^{(1)} \prec q^{(2)} \prec q^{(3)} \prec\cdots.
\]

The tower is analogous to iterated self-reference or 
meta-meta-circular interpreters in classical semantics—but with discrete curvature governing termination.

% ================================================================
\section{The Polyxan Löb Theorem (Full Strength)}
% ================================================================

Here we present the complete stabilisation theorem, the comparison to the 
λ-calculus Y combinator, and the full proof.

\begin{theorem}[Polyxan Löb Theorem]
\label{thm:polyx-loeb}
Let $G$ be any finite typed Polyxan hypergraph.  
Define the quine tower:
\[
G_0 = G,\qquad G_{n+1} = \mathcal{Q}(G_n).
\]
Then:
\[
G_3 \;\cong\; G_2 \;\cong\; G_1 \;\cong\; G_\infty.
\]
Thus the tower stabilises at exactly order $3$.  
Moreover,
\[
G_\infty
\]
is the initial $\mathcal{Q}$-coalgebra and terminal $\mathcal{Q}$-algebra over $G$.
\end{theorem}

\subsection*{Interpretation}

The Löb Theorem asserts that Polyxan self-reference is bounded:  
three layers of rewriting one’s own rewriting rules suffice to reach the universal fixed point.

This is the discrete analogue of the RSVP renormalisation group reaching  
the conformal fixed point $\mathcal{L}^{\ast}$ in three coarse-graining steps.

% ---------------------------------------------------------------
\subsection{Side-by-Side Comparison with the λ-Calculus Y Combinator}
% ---------------------------------------------------------------

\begin{center}
\begin{tabular}{l|l|l}
\textbf{Property} &
\textbf{λ-calculus (Y combinator)} &
\textbf{Polyxan $\mathcal{Q}^3$ fixed point} \\
\hline
Self-reference depth &
Infinite &
Finite, exactly 3 layers \\
Typing &
Requires untyped or impredicative types &
Fully typed, no impredicativity \\
Normalisation &
Non-terminating in untyped λ &
Always terminating (curvature descent) \\
Fixed point &
$Y M = M(YM)$ &
$G_\infty = \mathcal{Q}(G_\infty)$ \\
Universality &
Universal for recursion &
Universal for self-rewriting \\
Stability &
Divergent without evaluation strategy &
Stabilises at $3$ steps, independent of strategy \\
RSVP analogue &
None (λ lacks entropy) &
Exact analogue of $\mathcal{L}^\ast$ fixed point
\end{tabular}
\end{center}

This is the deepest structural divergence between classical recursion theory and Polyxan semantics.

% ---------------------------------------------------------------
\subsection{Full Proof of Theorem~\ref{thm:polyx-loeb}}
% ---------------------------------------------------------------

\paragraph{Step 1 — Curvature collapse.}
Each quine completion satisfies:
\[
\mathcal{H}_0[G_{n+1}] = \mathcal{H}_0[G_n] - 1,
\]
unless $G_n$ is already a media quine.  
Thus:
\[
\mathcal{H}_0[G_0] > \mathcal{H}_0[G_1] > \mathcal{H}_0[G_2] > \mathcal{H}_0[G_3] \ge 0.
\]

\paragraph{Step 2 — Stabilisation.}
Suppose $G_3\not\cong G_2$.  
Then $\mathcal{Q}(G_2)$ adjoins a new atom $q_{G_2}$.  
But $G_2$ already contains a self-rewriting atom encoding the full EHR system of
$G_1$ and $G_0$.  
Thus $q_{G_2}$ is isomorphic to an existing atom, forcing:
\[
G_3 \cong G_2.
\]
Repeating yields:
\[
G_2 \cong G_1.
\]

\paragraph{Step 3 — $G_\infty$ is a media quine.}
Self-fixpoint:
\[
\mathcal{Q}(G_\infty) \cong G_\infty.
\]
Thus $G_\infty$ is a media quine (Chapter~\ref{ch:hypergraph-fibrations}).

\paragraph{Step 4 — Initial/terminal structure.}
$\mathcal{Q}$ is left adjoint to $U$, so $G_\infty$ is the free quine-coalgebra on
$G$.  
By uniqueness of normal forms, it is also terminal among its quine-enrichments.

This completes the proof.

% ---------------------------------------------------------------
\subsection{Corollary: The Three-Storey Semantic Universe}
% ---------------------------------------------------------------

\begin{corollary}[Three-Storey Universe]
Every finite Polyxan hypergraph reaches complete self-description after 
exactly three layers:
\[
\text{``rewrite''},\quad
\text{``rewrite the rewriting rules''},\quad
\text{``rewrite the rule-of-rule-rewriting''}.
\]
No higher metalanguage is needed.
\end{corollary}

This is the discrete analogue of the RSVP RG flow stabilising at the conformal
fixed point $\mathcal{L}^\ast$ in three steps.

% ================================================================
\section{The Universal Media Quine}
% ================================================================

Define the higher-quine monad:
\[
\mathbb{Q}(G)=\mathcal{Q}^3(G).
\]

\begin{theorem}[Universal Media Quine]
The universal media quine $\mathfrak{Q}$ is the initial algebra of the
higher-quine monad:
\[
\alpha : \mathbb{Q}(\mathfrak{Q}) \to \mathfrak{Q},
\]
and for any quine-stable $G$ there exists a unique algebra morphism 
$\mathfrak{Q}\to G$.
\end{theorem}

This is the precise discrete analogue of the semantic conformal vacuum 
$\mathcal{L}^{\ast}$.

% ================================================================
\section{Summary}
% ================================================================

In this chapter we have:

\begin{itemize}
\item defined self-rewriting atoms and quine-enriched hypergraphs,
\item constructed the quine operator and iterated quine tower,
\item proved the Quine Fixed-Point Theorem,
\item defined higher-order quines and established the Polyxan Löb Theorem,
\item compared Polyxan self-reference with the λ-calculus Y combinator,
\item proved the three-storey universe corollary,
\item identified the universal media quine as the initial algebra of the
higher-quine monad.
\end{itemize}

This completes the self-referential backbone of Polyxan algebra and prepares for
the global RSVP–Polyxan coupling developed in Part~III.
